\Needspace{13\baselineskip}
\section{R5: General recursion with explicit obligations}\label{sec:recursion}

\Q{R5.Q1}{A capability interpreter and Lean's recursors}{Will a synthesis-time interpreter agree definitionally with \code{brecOn} or \code{WellFounded.fix}, and how should evaluation and publication work?}

\answer Not automatically. Fix a domain $D$, a family $B:D\to\mathsf{Sort}\ u$, a relation $R$ on $D$, and a proof that $R$ is well founded. A recursive body should have the shape
\begin{equation}\label{eq:capability}
 F:\prod_{x:D}\left(\prod_{y:D}R(y,x)\to B(y)\right)\to B(x).
\end{equation}
At the current input $x$, the inner function is a capability to call recursively only at a smaller $y$, with an explicit decrease proof. Its dependent result is $B(y)$, not always $B(x)$ or a fixed output type. Synthesis may leave the decrease proof as an obligation; it must not give the body an unrestricted self-reference and hope the final gate discovers termination later.

The well-founded fixpoint $f$ satisfies the equation
\[
 f(x)=F\bigl(x,\lambda y\ h.\ f(y)\bigr).
\]
Whether a particular equation is definitional depends on the representation, recursor, proof terms, and reduction policy. Lean's general well-founded equations are normally reasoned about propositionally. There is also a version-sensitive exception worth knowing: Lean~4.27 introduced a reducible fuel-based implementation for certain natural-number-measure definitions. One should neither promise universal definitional agreement nor repeat the outdated blanket claim that no well-founded definition can reduce in the kernel.\cite{leanrecursion,lean427}

An independent synthesis interpreter needs a simulation theorem. For a fuel interpreter, a useful statement is: if evaluation at $x$ returns $v$ with enough fuel, then $v=f(x)$; and, when a finite recursion-height bound is available, evaluation above that bound succeeds. Prove it by induction on the evaluation derivation or by well-founded induction on $R$. Fuel exhaustion proves nothing about the contract. A type-correct interpreter is not, merely by having the same input/output types, an implementation of the same function.

\subsubsection*{A practical publication path already suggested by Lean itself}

The 4.34.0 standard library contains a directly relevant precedent. It defines \code{Nat.recCompiled}, proves \code{Nat.rec_eq_recCompiled}, and marks the equality with \code{@[csimp]}. Its comment explicitly identifies the compiler's recursor-support limitation as the reason. This is a proof-backed compilation rewrite, not an unproved \code{implemented_by} substitution.\cite{leannat}

Leant2 should investigate this route for its recursor publication problem before developing a universal custom interpreter. Synthesize the semantic recursor term; construct or reuse an ordinary recursive definition; prove their equality, or directly reprove the contract for the published definition; then use an appropriate proof-backed compilation rewrite or publish the checked ordinary definition itself. An equality proof may use axioms such as those behind function extensionality, so the chosen policy must be audited rather than assumed unchanged. A pointwise or direct contract proof can sometimes avoid an unnecessarily strong equality obligation.

This still separates logic from runtime. A proof that two Lean definitions are equal certifies their logical relationship. Executable behavior additionally relies on the compiler/runtime and the absence of unaccounted-for implementation substitutions. Report successful export and compilation separately, and run representative observations on the exported artifact. Do not label a kernel-accepted but noncompiling recursor expression as an executable synthesized program.

\Q{R5.Q2}{What top-down angelic recursion must preserve}{Can capabilities and angelic evaluation be combined with top-down synthesis without losing the useful composition properties of Burst?}

\answer They can, but the missing ingredient is a shared constraint representation, not a different spelling of recursion. A partial recursive-equation IR should record its formal parameters, guards, branches, recursive call sites, holes, and pending termination proofs. Every occurrence of a hole points to the same unknown function body. Different example executions constrain that body in different environments.

An angelic return value summarizes an assumption on a call; it is not an implementation. When the submitted examples specify a smaller call, its observed result can constrain the oracle. Otherwise use a symbolic value and accumulate relations imposed by the enclosing context. At realization time, the same synthesized body must satisfy all these relations and every decrease obligation. A local combination that works with mutually inconsistent recursive answers must be rejected or refined before publication.

The main advantage to retain from a bottom-up system is representation of many candidate subprograms through a shared summary. A purely depth-first top-down loop can repeatedly rediscover the same realizations and lose that sharing. Add a typed library of closed candidates and compact version spaces for a supported first-order fragment. Conversely, top-down hole contexts retain the local type and induction information that an excessively erased bottom-up representation would lose. The hybrid should exchange typed recipes and constraints rather than mutually incompatible internal languages.

Smyth is relevant here precisely because it addresses interdependent goals without trace-complete examples. Burst offers an angelic strengthening perspective. The 2026 Cataclyst work provides a further comparison point for recursive sketches and mixed-quantifier first-order properties, with counterexample-derived syntactic pruning; it is not an off-the-shelf dependent Lean solution.\cite{smyth,burst,cataclyst} The Leant2 evaluation should compare these architectural ideas on matched specification regimes instead of transferring their aggregate success rates to a different language and workload.

\Q{R5.Q3}{A small recursion basis and evidence for coverage}{Is there a canonical small collection of recursion schemes that covers the intended tasks, and what coverage can be claimed?}

\answer There is no unique useful ``canonical basis'' independent of the language and cost model. A small implementation basis is still sensible: structural recursion over one selected inductive argument with generalized parameters; product/function accumulators; and well-founded recursion through \cref{eq:capability}. Helper introduction should initially be bounded to one explicitly synthesized helper, with mutually recursive and more elaborate schemes introduced in later size grades.

The structural tier is more capable than it may appear. List zip can recurse on the first list while the motive retains the second as a parameter, performing a case analysis on it in the step. Indexed zip requires the corresponding dependent motive. A tree fold can combine the results of both recursive children. Fibonacci can use a pair carrier instead of unrestricted course-of-values recursion. These examples justify the proposed ordering, not a measured coverage percentage.

The well-founded tier handles tasks whose direct equations do not follow one structural subterm. Merge can use the sum of the two list lengths. Euclidean gcd can use the second natural argument, with the nonzero branch proving that the remainder is smaller. Filtering recursive arguments in quicksort requires a length argument and appropriate inequalities. These are not supplied by an arbitrary call capability for free: guard facts, arithmetic lemmas, and contract proof search must solve them.

Coverage should be measured by ablation on three complementary groups: elementary list/natural/tree exercises; the unresolved Church and indexed cases already in the repository; and adversarial variants that alter index expressions, parameter order, or termination measures without changing the essential algorithm. Count separately tasks expressible in the grammar, tasks for which a witness was found, tasks with a checked contract, and tasks successfully exported and executed. Give both logical-work and wall-clock curves. No experiment in this review establishes a coverage percentage for the proposed basis.

A useful diagnostic comes from a reference solution withheld from synthesis. Check whether that solution can be represented within each tier's grammar and bound. If it cannot, a search failure is a coverage limitation. If it can, the failure is attributable to search order, checking cost, or implementation. Keeping these cases separate prevents a ``faster search'' claim from hiding the removal of difficult but intended programs.
